%********************************************************************
% Overall conclusion and future work
%********************************************************************
\chapter{Overall Conclusion and Future Work}

\section{Overall Conclusion}
The promises of scaling laws, where scaling up model size gives rise to new capabilities, 
are colliding with a harsh reality: training these large models and running inference 
is costly. The whole research question is how to scale language models efficiently,
both for training and inference. It is fortunate that pre-trained LLMs exhibit
redundancy at different levels, that can be leveraged to make them efficient
at post-training or pre-training. We show that one example of 
such redundancy is expressed in the fact that pre-trained Transformer's exhibit 
low-rank structure in their activations.
Our contributions exploit this observation to make LLMs efficient mainly in two
directions: \textbf{model compression} using activations and \textbf{activation sparsity}.

For post-training \textbf{model compression}, we propose \textbf{Lillama}, a method that locally distills the
activations of Transformer layers into low-rank weights. This local optimization makes
the approach more efficient than prior work: it enables compressing, on a single GPU,
LLMs that would not otherwise fit on one. To achieve speedup, the method does not require
specialized GPU kernels, unlike weight-sparsity approaches such as Wanda or SparseGPT.
The approach is also sample-efficient, requiring only 13M tokens, compared to methods
like Sheared-Llama, which uses 50B tokens for compression, or Minitron, which uses over
100B. Sample efficiency is an important property in post-training compression, as it
directly affects the fairness of the approach---for instance, its viability for
under-represented languages. This is why we specifically studied Whisper compression in
a low-resource setting, using Bambara as a case study. In this work, \textbf{BaldWhisper}, we
again distill input/output embedding activations into low-rank weights, but push
compression further by reducing the number of layers. To preserve the method's sample
efficiency, we do this by merging layers rather than removing them,
to stay close to the original model.
The resulting model is 2x faster while maintaining most of the performance.

For \textbf{activation sparsity}, we first study the eigenstructure of \textbf{Matryoshka
embedding LLMs} in multilingual and multimodal representation learning. We show
that Matryoshka embeddings, viewed as a form of structured activation sparsity,
leave room for more effective use of the representation space: most of the
activation energy is concentrated in relatively few principal components.
The fixed slicing opertor applied on the embedding vectors 
does not preserve all the information needed for downstream tasks while
introducing redundancy. We make the same observation for dense MLPs in
autoregressive LLMs: their intermediate activations occupy a subspace
substantially smaller than the allocated dimension. Sparse MoEs, however,
project each token into a smaller subspaces of the available intermediate dimension, 
and we showed that those small subspaes are used more fully. 
These observations motivate an approach where the MLP has a large intermediate
dimension but each token select only the needed intermediate coordinates adaptively 
rather than activating the whole available space.
We introduce \textbf{ZeroNeuron (ZoN)}, an input-dependent
gating mechanism that learns which coordinates to activate while constraining. 
With approximately 600M active parameters per
token, the 1.2B ZoN model matches the aggregate performance of its dense
counterpart and outperforms a dense model of comparable active parameters.
These results show that model capacity can be preserved while reducing
the parameters activated for each input, and the such sparsity
can be exploited by GPU Kernels to accelerate training and inference.

\section{Future Work}
We believe in sparsity, and we expect the most performant neural systems
of the future to be sparsely activated. The remaining question is how
this sparsity should be designed. We do not believe that TopK-based
mixtures of experts or sparse memory layers are the right long-term
answer: there is no reason to impose the same activation budget on
every token, regardless of the task or its difficulty. Moreover, under
current hardware constraints, distributed TopK routing incurs substantial
Cross-GPU expert communication is one of the biggest challenges in
scaling fine-grained sparsity, and approaches such as LatentMoE are
trying to address it. Maybe this will not be a problem
in few years with the next generation of hardware where a single GPU 
with TB of vRAM that can train multi-trillion parameters LLMs on a single GPU. 
But beyond the engineering issues, TopK-based
routing also faces the stability problems discussed earlier, including
dead experts and expert collapse. Rather than putting patches everywhere
to keep TopK-based sparsity working, we believe it can be beneficial to fundamentally
revisit how sparsity is learned. We bilieve that such approaches must be adaptive, communication-efficient,
and stable by design. ZoN is our first step in this direction. We plan
to investigate its scalability through larger models and longer training
runs, while developing GPU kernels that translate its learned sparsity
into actual computational savings.

The conclusion of this thesis is therefore the beginning of many things
to come, rather than the end of this research. We believe that learning
when, where, and how much to compute will be central to building more
capable and efficient neural systems.